\section{Pole structure and completeness}\label{sec:completeness}

This section proves Theorem~\ref{thm:complete} below: for odd \(p\),
and for even polynomials \(P\), every nonconstant meromorphic solution
of \eqref{eq:main} is rational, rational in an exponential, or
elliptic, and its poles form a single orbit under its periods.
The proof has three steps. We first compute the Laurent expansion at
a pole and show that for odd \(p\) it is unique. For even \(P\) it is
not: one coefficient is left free, and we show that the energy, a
first integral of \eqref{eq:main}, fixes it, so that Eremenko's
theorem applies on each energy level. Finally we show that any two
poles of one solution differ by a period.

Throughout, a meromorphic function is single-valued and meromorphic
on the whole complex plane. Let \(\Wclass\) consist of rational
functions of \(z\), rational functions of \(e^{\kappa z}\) with
\(\kappa\ne0\), and elliptic functions.
The period group of a nonconstant meromorphic function is its
\emph{full} group of periods: trivial for a rational function,
infinite cyclic for a rational function of \(e^{\kappa z}\), where
we choose \(\kappa\) so that \(2\pi i/\kappa\) generates it, and the
period lattice for an elliptic function.
For the general analytic background, see \citet{Eremenko1982}.

The theorem collects the analytic input to the enumeration.
Its odd-order case is contained in
\citet[Theorem~1]{YuanLiQi2014}.
For the even-order case, the first-integral mechanism is that of
\citet[Lemmas~3--4]{EremenkoLiaoNg2009}; here it is written uniformly
for every even polynomial \(P\). We give the proofs in full because
Sections~\ref{sec:systems} and~\ref{sec:counting} use the Laurent
recurrence and the energy dependence in explicit form.

\begin{theorem}[Completeness in the covered operator classes]
\label{thm:complete}
Assume either that \(p\) is odd or that \(P\) is an even polynomial.
Every nonconstant meromorphic solution of \eqref{eq:main} belongs to
\(\Wclass\). Every pole has order \(p\), and the pole set is one coset
of the full period group. Consequently there is one pole in the
rational case, one pole modulo the primitive period in the simply
periodic case, and one pole on the full elliptic period torus.

For fixed \(P\) of odd degree, there is at most one nonconstant solution
with its pole at zero. For even \(P\), the same statement holds after
fixing the conserved energy. The only entire solutions are
\(v=0\) and \(v=-2a_0\).
\end{theorem}

An even polynomial is one in which every odd power of \(\rho\) has
coefficient zero, so that \(P(D)\) contains only even derivatives.
The hypothesis does not cover an arbitrary operator of even degree,
and it cannot: Section~\ref{sec:open} gives a mixed second-order
operator with a meromorphic solution outside \(\Wclass\).

\subsection{Pole order and the Laurent recurrence}

Write a pole expansion at zero as
\begin{equation}\label{eq:laurent}
 v(z)=\sum_{n\ge0}c_nz^{n-p}.
\end{equation}
The pole order is forced by dominant balance. A pole of order \(r\)
makes \(P(D)v\) of order \(r+p\) and \(v^2\) of order \(2r\); the two
can cancel only if \(r=p\). Comparing the leading coefficients gives
\begin{equation}\label{eq:leadingpole}
 c_0=c_\ast:=-2(-1)^p R_p,\qquad
 R_p=\frac{(2p-1)!}{(p-1)!}.
\end{equation}
For the remaining coefficients, set
\begin{equation}\label{eq:indicial}
 I_p(n)=\prod_{j=p}^{2p-1}(n-j)+c_\ast.
\end{equation}
The product is the factor produced by \(D^p\) acting on \(z^{n-p}\),
and \(c_\ast\) is the factor produced by the cross term of
\(c_nz^{n-p}\) with the leading pole in \(v^2/2\).
With \(c_n=0\) for \(n<0\), comparison of the coefficients of
\(z^{n-2p}\) in \eqref{eq:main} gives
\begin{equation}\label{eq:recurrence}
 I_p(n)c_n=
 -\sum_{j=0}^{p-1}a_j
       \fall{(n-2p+j)}{j}\,c_{n-p+j}
 -\frac12\sum_{k=1}^{n-1}c_kc_{n-k},\qquad n\ge1.
\end{equation}
The underline denotes a falling factorial.
The right side contains only coefficients with index less than \(n\),
so \eqref{eq:recurrence} determines \(c_n\) whenever \(I_p(n)\ne0\).
The positive integer zeros of \(I_p\) are the \emph{resonances} of the
recurrence: at such an index the recurrence leaves \(c_n\)
undetermined if its right side vanishes, and has no solution
otherwise.

If \(p\) is odd, then \(c_\ast=2R_p\), and there is no positive
integer resonance.
For integers \(1\le n<p\), the product in \eqref{eq:indicial}
is negative with absolute value less than \(R_p\).
For \(p\le n\le2p-1\) it is zero, and for \(n\ge2p\) it is positive.
Hence \(I_p(n)\ne0\) for every positive integer \(n\), and there is
exactly one formal Laurent expansion at the marked pole.

To compare directly with \citet{YuanLiQi2014}, put \(w=-v/2\).
Then \(w^{(p)}+\sum_{j<p}a_jw^{(j)}=w^2\).
Their pole order is \(p\), their number of principal parts is one,
and each lower linear monomial has weighted pole order
\(p+j<2p\). Their degree and dominant-term hypotheses therefore hold.
We nevertheless apply Eremenko's theorem directly, since the same
theorem will be applied to an energy level in the even case.

Theorem~2 of \citet{Eremenko2006} states that an autonomous polynomial
differential equation with finitely many possible formal pole
expansions and a single monomial of highest total degree has all its
meromorphic solutions in \(\Wclass\).
In \eqref{eq:main} the only monomial of degree two is \(v^2/2\), and
for odd \(p\) we have just seen that there is exactly one formal pole
expansion. Both hypotheses hold at every odd order, for every
coefficient vector \(a\).

\subsection{The energy argument for an even operator}

Let \(p=2m\) and
\begin{equation}\label{eq:evenP}
 P(D)=D^{2m}+\sum_{j=0}^{m-1}b_jD^{2j}.
\end{equation}
Now \(c_\ast=-2R_p\), and there is exactly one positive integer
resonance.
Indeed the product in \eqref{eq:indicial} is at most \(R_p\)
for \(0\le n<p\), is zero for \(p\le n\le2p-1\), and is
strictly increasing for \(n\ge2p\); its value at \(3p\) is
\(2R_p=-c_\ast\). The only positive integer zero of \(I_p\) is
therefore \(n=3p\), and the possible free Laurent coefficient is
\[
 h=c_{3p}=[z^{2p}]v.
\]
If the recurrence is compatible at \(n=3p\), the equation has a
one-parameter family of formal expansions at the pole, one for each
value of \(h\), and Eremenko's theorem does not apply to
\eqref{eq:main} itself. The remedy is a first integral whose value
determines \(h\).

For \(j\ge1\), define the quadratic differential expression
\begin{equation}\label{eq:energyterms}
 \mathcal B_j(v)=
 \sum_{\ell=0}^{j-2}(-1)^\ell
       v^{(\ell+1)}v^{(2j-1-\ell)}
 +\frac{(-1)^{j-1}}2\bigl(v^{(j)}\bigr)^2,
\end{equation}
where the sum is empty when \(j=1\).
Termwise differentiation gives
\(D\mathcal B_j=v'v^{(2j)}\): every even derivative, multiplied by
\(v'\), is an exact derivative. Since \(P\) contains only even
derivatives, it follows that
\begin{equation}\label{eq:energy}
 \mathcal H_P(v)=
 \mathcal B_m(v)+\sum_{j=1}^{m-1}b_j\mathcal B_j(v)
 +\frac{b_0}{2}v^2+\frac16v^3
\end{equation}
satisfies
\begin{equation}\label{eq:energyidentity}
 D\mathcal H_P=v'E_P(v).
\end{equation}
In particular \(\mathcal H_P\), the energy, is constant on each
solution.
The case with no intermediate derivatives is the first integral
used by \citet{EremenkoLiaoNg2009}.
The familiar Kawahara first integral is the case \(m=2\)
\citep{DeminaKudryashovKawahara2010,YeEtAl2022};
the corresponding order-six construction appears in
\citet{Dang2023Even}.

The energy identity has two consequences for the Laurent expansion.
The first is that the resonance is always compatible.
To see this, construct \eqref{eq:laurent} through index \(3p-1\).
The first possible residual of \eqref{eq:main} is
\(\Omega_pz^p\).
Multiplication by \(v'\) gives residue \(-p c_\ast\Omega_p\).
By \eqref{eq:energyidentity} this is the residue of a derivative
of a Laurent series, and is therefore zero. Hence \(\Omega_p=0\),
and \(h\) may be chosen arbitrarily.

The second consequence is that the energy determines \(h\).
On the full formal solution the constant energy has the form
\begin{equation}\label{eq:energyh}
 \mathcal H_P=\Lambda_p h+\eta_p(P),\qquad
 \Lambda_p=-p c_\ast I_p'(3p)\ne0.
\end{equation}
To find the coefficient, linearize the highest-weight part of
\eqref{eq:energyidentity} about \(c_\ast z^{-p}\), with perturbation
\(z^{n-p}\).
If the corresponding energy term is \(L(n)z^{n-3p}\), then
\[
 (n-3p)L(n)=-p c_\ast I_p(n).
\]
Taking \(n\to3p\) gives the coefficient in \eqref{eq:energyh}.
Lower derivative terms have different powers of \(z\) and cannot
contribute to this coefficient. Moreover,
\[
 I_p'(3p)=2R_p\sum_{j=p+1}^{2p}\frac1j>0.
\]
Fixing the energy thus fixes \(h\), and all subsequent Laurent
coefficients follow from \eqref{eq:recurrence}.
This is the resonance-removal step in
\citet[Lemma~4]{EremenkoLiaoNg2009}, with an explicit coefficient
for the present operator family; see also
\citet[Section~3.2]{ConteNgWu2018} and
\citet[Section~5]{ConteMusetteNgWu2025}.

We can now apply Eremenko's theorem, not to \eqref{eq:main} but to
the level equation \(\mathcal H_P(v)=H\) with \(P\) and \(H\) fixed.
This equation is autonomous and polynomial, and its unique monomial of
highest degree is \(v^3/6\).
It has exactly one formal pole expansion: differentiating it and
dividing by the nonzero Laurent series \(v'\) returns \(E_P(v)=0\),
so every expansion belongs to the family above, and \eqref{eq:energyh}
selects the value of \(h\).
Every meromorphic solution of \eqref{eq:main} lies on one energy level
and therefore belongs to \(\Wclass\).

\subsection{The period group and the entire solutions}

It remains to show that the poles of one solution form a single orbit
under its periods.
For either covered case, take two poles \(z_1,z_2\) of the same
solution. The functions \(v(z+z_1)\) and \(v(z+z_2)\) satisfy the
same equation and, in the even case, have the same energy.
Their Laurent series at zero coincide, so the identity theorem gives
\[
 v(z+z_1)=v(z+z_2).
\]
Every difference of poles is a period. Conversely, every period
translates a pole to a pole. The period group is discrete: a sequence
of nonzero periods tending to zero would force \(v'\equiv0\).
The pole set is therefore one coset of that group.
The same argument proves the uniqueness statements of the theorem:
two nonconstant solutions with the same operator, the same energy in
the even case, and a pole at zero have the same Laurent series there.

Finally, an entire solution in \(\Wclass\) must be constant.
In the rational case it is a polynomial, and the
highest degree of its square cannot be cancelled in \eqref{eq:main}
unless the polynomial is constant.
In the rational-exponential case it is a Laurent polynomial in
\(t=e^{\kappa z}\); the operator \(P(\kappa t\partial_t)\) preserves
each exponent, whereas squaring doubles an extreme nonzero exponent.
An entire elliptic function is constant by compactness.
The constants solve \(a_0v+v^2/2=0\), which gives \(v=0\) and
\(v=-2a_0\).
This completes the proof of Theorem~\ref{thm:complete}.

The finite systems of the next section are derived from the one-pole
structure alone. The analytic facts of this section are used only to
show that, in the covered cases, those systems miss no nonconstant
meromorphic solution. The systems themselves never require the convergence of a
formal series.
